\par\Needspace{12\baselineskip}\bigskip
\pdfbookmark[2]{Darstellung, Zweiter Abschnitt}{pallas-II-section2}
\begin{center}
{\large\bfseries Zweiter Abschnitt.\par}
\medskip
\textit{Momentane Änderungen der Elemente infolge der Störungen.}
\end{center}
\pdfbookmark[3]{Darstellung, 7}{pallas-II-article7}
\begin{center}7.\end{center}
\noindent\hspace*{\parindent}Die eben erläuterte elliptische Bewegung setzt voraus, dass die gegenseitigen Anziehungen der Sonne und des Planeten die einzigen wirkenden Kräfte sind. Kommen noch andere Kräfte in Betracht, so wird die Bewegung des Planeten um die Sonne verwickelter. Wie diese Bewegung aber auch beschaffen sei, man sieht, dass sich der Planet, wenn diese Kräfte auf einmal aufgehoben würden, von diesem Augenblick an auf einem Kegelschnitt bewegen würde. Dessen Elemente wären Funktionen der Größen, die für diesen Augenblick den Ort und die erlangte Geschwindigkeit des Planeten bestimmen, das heißt der sechs Größen \(x,y,z,\dfrac{\mathrm{d}x}{\mathrm{d}t}\mathpunct{\raisebox{.45ex}{,}}\dfrac{\mathrm{d}y}{\mathrm{d}t}\mathpunct{\raisebox{.45ex}{,}}\dfrac{\mathrm{d}z}{\mathrm{d}t}\mathpunct{\cdot}\) Wir beschränken uns auf den Fall, in dem dieser Kegelschnitt stets eine Ellipse ist, \SourceLeaf{456}und nennen sie die \textls[80]{oskulierende Ellipse}. Man sieht also, dass alle endlichen Gleichungen zwischen den Elementen und den Größen \(x,y,z,r,\dfrac{\mathrm{d}x}{\mathrm{d}t}\mathpunct{\raisebox{.45ex}{,}}\dfrac{\mathrm{d}y}{\mathrm{d}t}\mathpunct{\raisebox{.45ex}{,}}\dfrac{\mathrm{d}z}{\mathrm{d}t}\mathpunct{\raisebox{.45ex}{,}}\dfrac{\mathrm{d}r}{\mathrm{d}t}\mathpunct{\raisebox{.45ex}{,}}\) die wir für den Fall der rein elliptischen Bewegung entwickelt haben, oder genauer alle Gleichungen, die für diesen Fall gelten, auch für die oskulierende Ellipse gelten. Da die Störkräfte aber ununterbrochen wirken, ändert sich die oskulierende Ellipse von Augenblick zu Augenblick. Ihre Elemente sind veränderliche Größen. Gelingt es jedoch, für jeden Augenblick die Zahlenwerte aller Elemente der oskulierenden Ellipse anzugeben, so ist die Aufgabe, die Bewegung zu bestimmen, gelöst. Man berechnet dann mit denselben Formeln wie bei der ungestörten Bewegung den Ort des Planeten ebenso wie seine Geschwindigkeit und die Richtung seiner Bewegung. Diese Untersuchung führt auf zwei Aufgaben. Zunächst bestimmen wir die momentanen Änderungen, welche die einzelnen Elemente durch die Wirkung der Störkräfte erfahren. Danach sind diese Differenzialausdrücke zu integrieren, um für jeden Augenblick die Werte der veränderlichen Elemente zu erhalten.
